\section{Correctness and running time}

\begin{numbered}{Theorem}{9.1}{correctness of the actual finite algorithm}
Run the dynamic program of Section 2 with $m$-element slots and $K=K_m$.
It accepts the padded set $J$ exactly when $J$ has a full schedule, and
returns a full schedule whenever it accepts. Consequently it decides the
original deadline and constructs an optimum by testing deadlines.
\end{numbered}
\uses{Lemmas 2.1, 2.2, 2.3 and 2.10; Proposition 2.8; Theorem 8.5.}
\begin{proof}
\step{1}{the finite search decides the padded instance [Proposition 2.8; Lemmas 2.3 and 2.10; Theorem 8.5]}
Proposition 2.8 turns every accepting state path into a full schedule.
Theorem 8.5 supplies a hierarchy for every full-schedulable $J$, and
Lemma 2.10 maps its marks and child gaps to an accepting state path of the
same finite search. Lemma 2.3 proves both directions of padding and its
removal.

\step{2}{elementary deadlines are decided without large padding [Lemmas 2.1 and 2.3]}
For a nonempty instance reject $T<1$. For $T\geq n$, the topological
order of Lemma 2.1 schedules one job per slot and avoids padding to a
possibly large $mT$. In the remaining branch reject $n>mT$; otherwise
run the padded algorithm.

\step{3}{the deadline trials return an optimum [Steps 1--2; Lemma 2.2]}
Test $T=1,\ldots,n$ in increasing order to optimize. Deadline $n$ is feasible, and the first accepted deadline is the
minimum. Lemma 2.2 also covers the original convention allowing real
start times. Assign the at most $m$ original jobs in a slot to distinct
machines in increasing identifier order.
\end{proof}

The normalized schedules in the structural proof certify completeness;
they need not be computed by the finite algorithm.

\paragraph{Encoding and the start of the padded bound.}
An explicit input uses canonical self-delimiting binary integers for
$m,n$, a signed integer $T$, an explicit list of vertices $1,\ldots,n$,
and a list of distinct arcs with endpoints in that list. Its bit length
$L$ includes delimiters and length fields, so $n\leq L$.
The graph is a DAG; malformed inputs can also be rejected by the
preprocessing below. Arbitrary finite binary vertex names can first be
renamed by repeated input scans. Input reading and that renaming are
charged in $L$. The bound in $N$ begins with the parsed, canonically
numbered padded DAG on $J=[N]$, where $N=mT\geq m$.

\begin{numbered}{Proposition}{9.2}{explicit multitape time bound}
For $m\geq3$ and $K\geq1$, the padded dynamic program, including its
witness on every accepted instance, runs on one deterministic multitape
Turing machine in
\[
 A(m,K)N^{5mK+4m+1}\tag{9.1}
\]
steps, where one effective coefficient bound is
\[
 A(m,K)=2^{40}(m+1)^4(K+1)^4(K512^K)^5.
\]
The machine has a fixed finite number of tapes, independent of $m$ and
$K$.
\end{numbered}
\uses{Lemmas 2.1 and 2.7; Proposition 2.8; the state and transition definitions in Section 2.}
\begin{proof}
\step{1}{at most $B$ formulas and $S$ states suffice [Lemma 2.7; state definitions]}
Put
\[
 C=K512^K,\quad B=CN^{mK},\quad S=BN^m,
 \quad h=\lceil\log_2(N+1)\rceil,\quad q=64(m+1)(K+1).
 \tag{9.2}
\]
These are bounds; loops visit only actual records. Lemma 2.7 bounds the
formula count by $B$. Explicitly, the leaf alphabet has at most
$4+8(N+2)+20\binom Nm\leq64N^m$ symbols. For $\ell$ leaves there
are at most $4^\ell$ tree shapes and $2^{\ell-1}$ internal Boolean
labellings. Summing their product with $(64N^m)^\ell$ over
$1\leq\ell\leq K$ gives $B$.
For fixed $W$, a described set $F$ and an $m$-subset $Z$ determine
$L=W\cap F$ and $R=W\setminus(L\cup Z)$. Reject invalid candidates.
Thus at most $B\binom Nm\leq S$ candidate states remain, including
duplicates.

\step{2}{all records have explicit scan bounds [encoding; elementary tape operations]}
Use a fixed collection of tapes (32 suffice) for the input, matrices,
rank map, formula stack, files, copied records, counters and output.
All variable-length stacks are delimited records on these tapes.
Store the reachability matrix in row-major order. A unary distance
scratch word permits a complete seek to a matrix location in $O(N^2)$
steps, with the scratch and matrix heads advancing together; no binary
counter is incremented for each matrix bit.

Store formula and described-set records in files with record length at
most $qN$. A set record contains its $N$ bits, cardinality, ordinal and
acceptance bit. A current-state record contains its three $N$-bit parts,
ordinal, reached/initial/final flags and predecessor ordinal, also in
at most $qN$ bits. Indeed,
\[
 \lceil\log_2(S+1)\rceil
 \leq 2+10K+(mK+m)h\leq qh/8\leq qN/8.
\]
Formula codes use $O(K(m+1)h)$ bits, within the same record envelope.
Fetch or update an indexed state by scanning its file and comparing
ordinal fields. This costs $O(qSN)$ including rewinds and all index
bits. Acceptance queries compare their $N$-bit query against the set
file. Keep saved schedules in a \emph{separate} file, with at most $N$
identifiers per schedule and at most $qN^2$ bits per record. A combined
acceptance/witness lookup costs $O(qBN^2)$.

For explicit slack let $D=2^{20}q^2$. In the scan implementation just
described, each row below costs at most $D$ times its displayed
expression. Scan, copy, compare and rewind use a bounded number of tape
moves per symbol; the records include their counters. At most 64
primitive passes per operation, including the postorder stack scans,
fit this coefficient. No record access has unit cost.

\begin{center}\small
\begin{tabularx}{\linewidth}{@{}p{0.50\linewidth}X@{}}
\toprule
Subroutine, over the whole padded call & Step bound divided by $D$\\
\midrule
Adjacency, reachability and rank & $N^5$\\
Generate and evaluate formulas & $BN^3$\\
Pairwise deduplication & $B^3N$\\
Cardinalities and their ordering & $B^3N^2$\\
Enumerate and initialize states & $BS(BN^2+N^4)$\\
State-graph propagation & $3BS^4N$\\
Reconstruct all saved witnesses & $BSN^2+B^2N^3$\\
Write the selected witness & $N^3$\\
\bottomrule
\end{tabularx}
\end{center}

\step{3}{graph preprocessing costs at most $DN^5$ [Lemma 2.1; Step 2]}
At most $N^2$ arc insertions cost $O(N^4)$ by full matrix seeks.
Floyd--Warshall uses $N^3$ triples, each with a bounded number of
$O(N^2)$ seeks, hence $O(N^5)$ steps. Its invariant is reachability
through the already processed internal vertices. Repeated scans select
the least active vertex without an active predecessor, append it to the
rank order and remove it; all removals fit $O(N^5)$ as well. Lemma 2.1
justifies that rank order.

\step{4}{formula generation and evaluation cost at most $DBN^3$ [Lemma 2.7; Steps 1--2]}
Generate valid Dyck words of length $2(\ell-1)$ by depth-first
extension through valid prefixes; they encode the ordered full binary
trees with $\ell$ leaves. There are $O(K)$ prefix edges per output and
$O(K)$ steps per update, giving $O(K^2)$ delay. Enumerate their Boolean
labellings and leaf parameters by finite counters. A cone leaf needs at most $mN$
reachability queries, each with an $O(N^2)$ seek. At most $K$ leaves
therefore cost $O(mKN^3)$ per formula. Threshold, slot and complement
leaves cost less. Postorder evaluation combines $N$-bit vectors; even
scanning the entire $O(KN)$-bit stack at each node costs $O(K^2N)$.
These parameter factors fit $D$, giving $DBN^3$.

\step{5}{deduplication and ordering fit their table bounds [Step 2]}
For each of at most $B^2$ record pairs, full file seeks fetch and compare
the vectors and mark the later duplicate. Each pair costs $O(qBN)$,
giving $DB^3N$; compact the retained records afterwards. Count each
vector's one-bits and store its cardinality. Repeated selection scans,
even with full record retrieval for every compared pair, cost at most
$DB^3N^2$. Break equal-cardinality ties by ordinal.

\step{6}{state enumeration costs at most $DBS(BN^2+N^4)$ [state definitions; Steps 1--2]}
For each of $B$ sets and at most $S$ candidates, fetch $W,F$, construct
$L,Z,R$, and test containment, disjointness, the antichain condition and
compatibility. Considering each ordered job pair and seeking to its
relation separately costs $O(N^4)$. The set fetches and initial/final
acceptance queries cost at most $DBN^2$. This gives the fifth row.

\step{7}{propagation costs at most $3DBS^4N$ [transition definition; Steps 1--2 and 6]}
Initially mark the initial states reached. Make at most $S$ complete
passes over at most $S^2$ ordered state pairs. On the first successful
transition to an unreached state, record its reached flag and its
predecessor; leave predecessors unchanged subsequently. Fetching or
updating state records costs at most $DSN$. An acceptance query costs
at most $DBN^2$, and even rechecking compatibility costs at most $DN^4$.
For $m\geq3$,
\[
 BN^2\leq SN=BN^{m+1},\qquad N^4\leq SN.
\]
A pair test therefore costs at most $3DSN$. Every reachable vertex has
a simple path of fewer than $S$ edges. Induction on path length proves
that $S$ full passes find every reachable vertex. Over all $W$ the cost
is
\[
 B\cdot S\cdot S^2\cdot3DSN
   =3DBS^4N=3DB^5N^{4m+1}.\tag{9.3}
\]

\step{8}{stored witnesses fit the last two rows [Proposition 2.8; Steps 2 and 7]}
Every transition increases $|L|$ by at least $m$, so a predecessor path
contains at most $N/m$ states. Backtrack, reverse the path records, and
retrieve the already stored schedules for its first left part,
transition gaps and final right part. State retrieval costs at most
$DSN^2$ per $W$; the at most $N/m+1$ schedule lookups cost at most
$DBN^3$. Proposition 2.8 proves that these pieces and marked slots
partition $W$, so the concatenated schedule contains exactly $|W|$
identifiers. Store it once, retaining the just-constructed schedule on a work tape,
with no recursive recomputation of smaller witnesses. Over $B$ sets this
gives the seventh row. The full set $J$ is the unique set of cardinality
$N$, so its final witness is available on that tape. Translation through
the rank map and output fit the final $DN^3$ row.

\step{9}{the sum has exponent $5mK+4m+1$ [Steps 1--8]}
The state-enumeration row expands to
$B^3N^{m+2}+B^2N^{m+4}$, and the witness row to
$B^2N^{m+2}+B^2N^3$. Each has $B$-exponent at most five and
$N$-exponent at most $4m+1$, since $m\geq3$. The other rows satisfy
the same inequalities; in particular $5\leq4m+1$.
Thus their sum is at most $16DB^5N^{4m+1}$. Substituting $B$ gives
$16DC^5N^{5mK+4m+1}$, and
$16DC^5\leq A(m,K)$ by the definitions of $q,D,A$.
This proves the complete bound.
\end{proof}

\begin{numbered}{Corollary}{9.3}{the bounds of Theorem 1.1}
For $K_m=12m+25$, the padded exponent is
$e_m=60m^2+129m+1$. Deadline decision, its witness, and optimization
on the explicit encoding above take
\[
 O_m\bigl((L+2)^{60m^2+129m+2}\bigr)
\]
steps. One bound for the implicit coefficient is
\[
 G(m)=A(m,12m+25)m^{60m^2+129m+1}+2^{31}(m+1)^4.
\]
\end{numbered}
\uses{Lemma 2.3; Lemma 8.4 and Theorem 8.5; Theorem 9.1; Proposition 9.2.}
\begin{proof}
\step{1}{the padded exponent is $e_m=60m^2+129m+1$ [Lemma 8.4; Theorem 8.5; Proposition 9.2]}
Substitution gives
\[
 5m(12m+25)+4m+1=60m^2+125m+4m+1
 =60m^2+129m+1.
\]
\step{2}{parsing and elementary deadlines cost $P_m(L+2)^4$ [encoding; Lemma 2.1]}
Parse the input once. Repeated scans, endpoint comparisons, duplicate
checks, binary comparisons and canonical renaming fit
$P_m(L+2)^4$ steps with $P_m=2^{30}(m+1)^4$.
For $T\geq n$, produce a topological order directly from the original
arc list: at most $n$ removal stages and $n$ candidate vertices per
stage, with at most $O(L^2)$ steps to scan the arcs and active-vertex
vector for a candidate. This is $O(n^2L^2)\subseteq O(L^4)$.
The branch $T<1$ costs less. Both include their output cost.

\step{3}{all deadline trials add one power of $L$ [Lemma 2.3; Theorem 9.1; Proposition 9.2]}
In every padded branch, $N=mT\leq mn\leq mL$ by Lemma 2.3.
Theorem 9.1 tests at most $n\leq L$ deadlines. Proposition 9.2
bounds their total padded work by
\[
 A(m,K_m)n(mn)^{e_m}
 \leq A(m,K_m)m^{e_m}(L+2)^{e_m+1}.
\]
\step{4}{output and preprocessing fit the same bound [Steps 1--3; encoding]}
Writing the selected original schedule uses at most
$O(n(\log(n+1)+\log(m+1)))$ bits and fits the preprocessing envelope.
Since $e_m+1\geq929>4$, adding these costs gives $G(m)$ and the
stated exponent. The extra power beyond $e_m$ counts the deadline
trials. The padded bound begins after parsing, so a long original input
is never silently charged to $N$.
\end{proof}

\begin{center}\small
\begin{tabular}{@{}rrrrrr@{}}
\toprule
$m$ & union leaves & global leaves & relative leaves & $N$ exponent & $L+2$ exponent\\
\midrule
3&27&61&34&928&929\\
4&33&73&40&1477&1478\\
5&39&85&46&2146&2147\\
6&45&97&52&2935&2936\\
\bottomrule
\end{tabular}
\end{center}

\begin{numbered}{Proposition}{9.4}{the compact-isolated-jobs remark}
For fixed $m$, deadline decision is polynomial when the input specifies
$n$ in binary and lists only the arcs, allowing many isolated jobs.
An explicit schedule can require a number of output symbols
exponential in that compact input length.
\end{numbered}
\uses{Lemmas 2.1 and 2.3; Theorem 9.1; Proposition 9.2.}
\begin{proof}
\step{1}{the capacity test and endpoint count are valid [schedule and arc-list definitions]}
For a nonempty instance reject $T<1$ or $n>mT$. Let $d$ be the number
of distinct arc endpoints, found and renamed by input scans; then
$d\leq2|E|$.

\step{2}{deadlines $T\geq d$ are feasible after the capacity test [Lemma 2.1; Step 1]}
For $T\geq d$, Lemma 2.1 schedules the $d$ incident
vertices one per slot. There are $mT-d\geq n-d$ vacant positions, so
all isolated jobs can be placed injectively in them. This proves
feasibility without explicitly listing those jobs.

\step{3}{deadlines $T<d$ permit polynomial expansion [Step 1; Lemma 2.3; Theorem 9.1; Proposition 9.2]}
For $T<d$, the capacity check and the endpoint bound give
\[
 n\leq N=mT<md\leq2m|E|.
\]
The entire vertex set and its padding can therefore be expanded in
polynomial time in the compact input length for fixed $m$.
Lemmas 2.3 and Theorem 9.1 identify the answer of the expanded decision
computation with the original deadline answer, and Proposition 9.2
bounds that computation. All endpoint-label comparisons are charged
to the compact input length.

\step{4}{an explicit output can be exponentially long [compact and explicit encodings]}
For fixed $m$ take $n=2^k$ isolated jobs, $E=\varnothing$, and $T=n$.
The compact encoding has $O(k)$ bits, while every literal schedule lists
all $2^k$ jobs and therefore has at least $2^k$ output symbols. The
assertion here concerns decision. The empty instance is handled directly
under the convention that its makespan is zero.
\end{proof}
